\scanpage{122}
% archive foliation: 121
\inserted{$\operatorname{Conn}_{X/S}$ et de}
\[
H\backslash\mathcal P_{X/S}\longrightarrow\operatorname{Conn}_{X/S}
\]
s'étendent à $G_0(t,s)$ et à
\[
H\backslash\mathcal P_{X/S}\longleftrightarrow\operatorname{Conn}_{0\,X/S},
\]
les premières se déduisant des secondes de façon évidente, compte tenu de l'hom. can.
\[
\operatorname{Conn}_{0\,X/S}\longrightarrow\operatorname{Conn}_{X/S}
\]
compatible avec l'hom.
\[
\underline\omega_{X/S}^{\otimes2}\xrightarrow{12\cdot\mathrm{id}}\underline\omega_{X/S}^{\otimes2}
\]
sur les Groupes d'opérateurs. (On \deleted{posera} \inserted{écrira} parfois $c_0^X(t)=\frac1{12}c^X(t)$ \ldots)

On peut espérer que si $S$ de car. $p>0$ et si $\mathcal O'_X$ est le ss-faisceau de $\mathcal O_X$ \deleted{défini par} \inserted{\unsure{à qui s'appliquent}} les conditions $D_t^{(i)}f=0$ ($i\le3$)
\[
H_{X^{(p/S)}}\backslash\mathcal P_{X/S}\lhook\joinrel\longrightarrow\operatorname{Conn}_{0\,X/S}
\]
soit \uline{injectif}, i.e. qu'on a
\margin{$\mathcal O'_{X/S}$ est $\mathcal O_{X/S}^{(p)}$ si $p\ne2,3$, sinon c'est $\mathcal O_{X/S}^{(p^2)}$.}
\margin{\uline{OK}}

\uline{Proposition} \deleted{\unsure{Non vérifiée}} $A$ anneau de car. $p>0$, \inserted{local (plus gén. \unsure{avec} ${}_2\operatorname{Pic}(A)=0$)}, $\alpha\in A$, $B=A[T]/(T^p-\alpha)$, $t\in B$ la classe de $\deleted{B}T$ dans $B$ (donc $t^p=\alpha$), $\unsure{\forall|\exists}\ s\in B$, donc
\[
s=a_0+a_1t+\cdots+a_{p-1}t^{p-1}\qquad a_i\in A\ (0\le i\le p-1).
\]
\margin{car. $p\ne2,3$ (\unsure{hypothèse provisoire})}
Posons
\deleted{$D_t^{(j)}s=\sum_{j\le i\le p-1}a_i\binom ij t^{i-j}$},
\inserted{$D_t^1s:=(ds)/(dt)$, $D_t^{(i)}s:=\frac1{i!}(D_t^1)^is$ si $1\le i\le3$, et}
\[
\deleted{\unsure{G_0}}\ c_0(s)=-\frac{D_t^{(\unsure{2|3})}s}{D_t^{(1)}s}+\left(\frac{D_t^{(\unsure{2|3})}s}{D_t^{(1)}s}\right)^2\qquad\text{si }D_t^{(1)}s\in B^*.
\]
Supposons $c_0(s)=0$ i.e.
\[
(D^{(3)}s)^2=(D^{(1)}s)(D^{\unsure{(2)|(3)}}s).
\]
Alors $\exists a,b,c,d\in A$ \inserted{(uniques à \unsure{transformation} \uline{diagonale})} avec $ad-bc=1$, tels que
\[
s=\frac{at+b}{ct+d}\qquad\text{i.e.}\quad s(ct+d)=at+b
\]
